\documentclass[11pt]{article}

\usepackage[a4paper,margin=1in]{geometry}
\usepackage{hyperref}
\usepackage{enumitem}
\usepackage{titlesec}

\setlist[itemize]{noitemsep, topsep=2pt}

\titleformat{\section}
  {\large\bfseries}
  {}
  {0em}
  {}[\titlerule]

\begin{document}

%====================
% Header
%====================

\begin{center}
{\Large \textbf{Tai, Wei Hsuan's Curriculum Vitae}} \\
\vspace{4pt}
\href{mailto:b12508026@ntu.edu.tw}{b12508026@ntu.edu.tw} \quad
\href{https://github.com/W-X-Dai}{github.com/W-X-Dai}
\end{center}

\vspace{10pt}

%====================
% Education
%====================

\section*{Education}

\textbf{National Taiwan University} \hfill Current \\
B.S. in Biomedical Engineering

%====================
% Research Interests
%====================

\section*{Research Interests}

Machine Learning, Deep Learning, Computer Vision, Medical Image Analysis


\section*{Technical Skills}

\textbf{Programming:} C++ (Codeforces max rating: 1376), Python \\ \\
\textbf{Machine Learning:} PyTorch, Diffusion Models, Self-Supervised Learning \\ \\
\textbf{Database:} Schema Design, SQL, PostgreSQL, Redis \\ \\
\textbf{Web Development:} NestJS, Vue, JavaScript \\ \\
\textbf{Tools:} Linux(Ubuntu), Docker, Git


%====================
% Research Experience
%====================
    
\section*{Research \& Project Experience}

\textbf{Database Management System Final Project - Uniconnect} \hfill 2025 Fall
\begin{itemize}
    \item \href{https://github.com/W-X-Dai/DBMS_FinalProj}{https://github.com/W-X-Dai/DBMS\_FinalProj}
    \item This is the final project for the Database Management System course, where we developed a database management system called Uniconnect. The purpose of Uniconnect is to manage the resources such as internships or gpu servers provided by the university or the enterprises. Enterprises can also find students with specific skills through Uniconnect. I mainly responsible for the backend development of Uniconnect, which is implemented in NestJS. The database is implemented in PostgreSQL, and we also use Redis for caching. The frontend is implemented in Vue by my teammates, and I utilize Docker to deploy the application. (This project is co-worked with LLMs)
\end{itemize}
\textbf{Text Recovery with Denoisy Unet} \hfill 2025 Fall
\begin{itemize}
    \item \href{https://github.com/W-X-Dai/text_recovery}{https://github.com/W-X-Dai/text\_recovery}
    \item This is originally a homework of \textit{Fundamentals of Biomedical Image Processing} course, it originally expected us to utilize CLAHE to make the text more clear. But I think the performance of CLAHE is not good enough, so I implemented a Denoisy Unet to recover the text. The Denoisy Unet is trained on the synthetic dataset generated by adding noise to the original images. The Denoisy Unet can effectively recover the text from the noisy images, and it can also generalize well to the real-world noisy images.
\end{itemize}
\textbf{Fundamentals of Biomedical Image Processing Final Project} \hfill 2025 Fall
\begin{itemize}
    \item \href{https://bit.ly/4lwuwQy}{https://bit.ly/4lwuwQy}
    \item This is the final project for the Fundamentals of Biomedical Image Processing course, where we developed a deep learning model to segment the brain tumor from the MRI images. We utilized the BraTS dataset for training and evaluation. In this work, I mainly responsible for the distillation methods. I distilled from a well-trained nnUNet weight to a 2.5D MobileNet based Unet. The distilled model can achieve comparable performance to the original nnUNet, while it is much smaller and faster. 
\end{itemize}
%====================
% Technical Skills
%====================


\end{document}